%%%%%%%%%%%%%%%%%%%%%%%%%%%%%%%%%%%%%%%%%%%%%%%%%%%%%%%%%%%%%%%%%%%%%%%%%%%%%%%%
\section{Nanopub Onboarding}
{   
	\usebackgroundtemplate{
		\vbox to \paperheight{\vfil\hbox to \paperwidth{\hfil\includegraphics[width=\paperwidth]{humor/onboarding_cartoon.jpg}\hfil}\vfil}
	}
	\frame{
		\frametitle{Register with the Nanopub}
		\begin{mdframed}[tikzsetting={draw=white,fill=white,fill opacity=0.8,
				line width=0pt},backgroundcolor=none,leftmargin=0,
			rightmargin=150,innertopmargin=4pt,roundcorner=10pt]
			\tableofcontents[currentsection,sections={1-4},hideothersubsections]
		\end{mdframed}
		\vfill\vspace{18mm}\hfil \hfill{\tiny own sketch after "Dropping the Pilot" -- CC-BY-4.0}
	}
}

%%%%%%%%%%%%%%%%%%%%%%%%%%%%%%%%%%%%%%%%%%%%%%%%%%%%%%%%%%%%%%%%%%%%%%%%%%%%%%%% 
\frame{
  \frametitle{What is a Nanopub?}
  \begin{columns}[T]
    \column{0.3\textwidth}
       \centering
       \includegraphics[width=\linewidth]{nanopub_logo.png} \\
       Think of a Nanopub as a little knowledge container.
    \column{0.7\textwidth}
       \centering
       \includegraphics[width=\linewidth]{example_nanopub_dataset.png} \\
       This is an example Nanopub, a dataset description.
   \end{columns}
}
    
%%%%%%%%%%%%%%%%%%%%%%%%%%%%%%%%%%%%%%%%%%%%%%%%%%%%%%%%%%%%%%%%%%%%%%%%%%%%%%%% 
\frame{
  \frametitle{Register your account}
  \begin{columns}
      \begin{column}{0.5\textwidth}
          If you don't have an account yet, please register at
          \href{https://nanodash.knowledgepixels.com}{nanodash} using your ORCID iD.
      \end{column}
      \begin{column}{0.5\textwidth}
         \includegraphics[width=0.9\linewidth]{register.png}
      \end{column}
  \end{columns}


  
}

%%%%%%%%%%%%%%%%%%%%%%%%%%%%%%%%%%%%%%%%%%%%%%%%%%%%%%%%%%%%%%%%%%%%%%%%%%%%%%%%  
\frame{
  \frametitle{Interlude: What is an ORCID?}
  ORCID stands for \textbf{O}pen \textbf{R}esearcher and \textbf{C}ontributor \textbf{ID}.
  \vspace{1em}
  \begin{itemize}
    \item<+-> free, unique \& persistent ID
    \item<+-> show your research interest, funding, employment and \textbf{publications} (sometimes they add automagically with Crossref)
    \item<+-> useful in research to be found:
    \begin{itemize}
       \item<+-> for agencies to look you up during proposal reviews
       \item<+-> to find colleagues (and see how actively they are engaged and in which topics)
    \end{itemize}
  \end{itemize}
}
   
%%%%%%%%%%%%%%%%%%%%%%%%%%%%%%%%%%%%%%%%%%%%%%%%%%%%%%%%%%%%%%%%%%%%%%%%%%%%%%%%  
\frame{
  \frametitle{Interlude: What is an ORCID? II}
  \begin{columns}[T]
    \begin{column}{0.25\textwidth}
      \centering
      \includegraphics[width=\linewidth]{orcid_profile_links_keywords.png} \\
      Example screenshot with links and research interests.
    \end{column}
    \begin{column}{0.65\textwidth}
        \centering
        \includegraphics[width=\linewidth]{orcid_paper_examples.png} \\
        Example screenshot of papers linked in an ORCID profile.
    \end{column}
  \end{columns}
}

%%%%%%%%%%%%%%%%%%%%%%%%%%%%%%%%%%%%%%%%%%%%%%%%%%%%%%%%%%%%%%%%%%%%%%%%%%%%%%%%     
\frame{
  \frametitle{Authorise ORCID iD and enter your Nanodash account}
   \begin{columns}[T]
      \column{0.5\textwidth}
      \centering
      \includegraphics[width=0.9\linewidth]{authorise_orcid.png}
      \column{0.5\textwidth}
      \centering
      \includegraphics[width=0.9\linewidth]{enter_nanodash.png}
   \end{columns}
}

%%%%%%%%%%%%%%%%%%%%%%%%%%%%%%%%%%%%%%%%%%%%%%%%%%%%%%%%%%%%%%%%%%%%%%%%%%%%%%%%     
\frame{
  \frametitle{Introduction of an User}
  Click on the profile icon in the top right corner to access your profile page.\newline
  \vspace{1em}
  \centering
  \includegraphics[width=0.8\linewidth]{profile.png} \\
  \vspace{1em}
  Then click on the \textbf{“See Your Profile Details”} button to access your profile details.
}

%%%%%%%%%%%%%%%%%%%%%%%%%%%%%%%%%%%%%%%%%%%%%%%%%%%%%%%%%%%%%%%%%%%%%%%%%%%%%%%%        
\frame{
  \frametitle{Create Introduction}
  Inside the Profile Details page, click on the \textbf{“Create Introduction”} button to create your introduction.
  \vspace{1em}
  \centering
  \includegraphics[width=0.8\linewidth]{intro.png}
}
    
%%%%%%%%%%%%%%%%%%%%%%%%%%%%%%%%%%%%%%%%%%%%%%%%%%%%%%%%%%%%%%%%%%%%%%%%%%%%%%%%        
\frame{
  \frametitle{And then ...}
  Once inside the introduction page, you will see a \textbf{“Create Nanopublication”} form.
  \vspace{1em}
  \centering
  \includegraphics[width=0.8\linewidth]{intro_nanopub.png}
}
    
%%%%%%%%%%%%%%%%%%%%%%%%%%%%%%%%%%%%%%%%%%%%%%%%%%%%%%%%%%%%%%%%%%%%%%%%%%%%%%%%        
\frame{
  \frametitle{Complete your introduction}
  Under the \textbf{“Introducing a user”} section, add your name as you would like it to appear on your profile.
  \vspace{1em}
  \centering
  \includegraphics[width=0.8\linewidth]{complete_intro.png} \\
  \vspace{1em}
  Finally, after you are happy, click the \textbf{“Publish”} button to publish your introduction.
}

%%%%%%%%%%%%%%%%%%%%%%%%%%%%%%%%%%%%%%%%%%%%%%%%%%%%%%%%%%%%%%%%%%%%%%%%%%%%%%%%    
\frame{
  \frametitle{A published introduction ...}
  Once the introduction is published, you should now be able to view your name on it.
  \vspace{1em}
  \centering
  \includegraphics[width=0.8\linewidth]{name_profile.png}
}

%%%%%%%%%%%%%%%%%%%%%%%%%%%%%%%%%%%%%%%%%%%%%%%%%%%%%%%%%%%%%%%%%%%%%%%%%%%%%%%% 
\frame{
  \frametitle{View the latest Nanopubs}
  Your name will appear under the latest nanopublications on your profile page.
  \vspace{1em}
  \centering
  \includegraphics[width=0.8\linewidth]{latest_nano.png}
}

%%%%%%%%%%%%%%%%%%%%%%%%%%%%%%%%%%%%%%%%%%%%%%%%%%%%%%%%%%%%%%%%%%%%%%%%%%%%%%%%    
\frame{
  \frametitle{Add “view display” to your profile}
  To add display sections to your profile, click on the \textbf{“+ view display”} button; a new page for adding a display section will open.
  \begin{columns}[T]
     \column{0.5\textwidth}
     \centering
     \includegraphics[width=0.9\linewidth]{view_display.png}
     \column{0.5\textwidth}
     \centering
     \includegraphics[width=0.9\linewidth]{create_view_display.png}
  \end{columns}
}
    
%%%%%%%%%%%%%%%%%%%%%%%%%%%%%%%%%%%%%%%%%%%%%%%%%%%%%%%%%%%%%%%%%%%%%%%%%%%%%%%%    
\frame{
  \frametitle{Approval of an User}
  Next on our list: approval of a user. Users can be approved by approved users.
  \vspace{1em}
  \centering
  \includegraphics[width=0.8\linewidth]{approval_with_highlighting.png} \\
  \vspace{1em}
  \textbf{Task:} approve your neighbour(s)!
}
    
%%%%%%%%%%%%%%%%%%%%%%%%%%%%%%%%%%%%%%%%%%%%%%%%%%%%%%%%%%%%%%%%%%%%%%%%%%%%%%%%    
\frame{
  \frametitle{Publish from a Template}
  \begin{columns}[T]
            \column{0.3\textwidth}
            \centering
            \includegraphics[width=\linewidth]{publish_menu.png} \\
            Select \textbf{“Publish”} from the menu.
            \column{0.7\textwidth}
            \centering
            \includegraphics[width=\linewidth]{publish_templates.png} \\
            Then you \textbf{can} choose from the list of templates. (We are going to publish together. Please wait.)
        \end{columns}
}

%%%%%%%%%%%%%%%%%%%%%%%%%%%%%%%%%%%%%%%%%%%%%%%%%%%%%%%%%%%%%%%%%%%%%%%%%%%%%%%%    
    \frame{
        \frametitle{Your first Nanopub!}
        Let’s select an easy template: search for \textbf{“Announcing a paper I have read”}.
        \vspace{1em}
        \centering
        \includegraphics[width=0.8\linewidth]{publish_paper_read.png} \\
        \vspace{1em}
        Select the template and enter the DOI of a recently read paper.
    }

%%%%%%%%%%%%%%%%%%%%%%%%%%%%%%%%%%%%%%%%%%%%%%%%%%%%%%%%%%%%%%%%%%%%%%%%%%%%%%%%    
    \frame{
        \frametitle{Modifying a Template}
        There are many templates for different purposes and disciplines. \textit{But}, sometimes you need a new template \dots
        \vspace{1em}
        \begin{callouttip}
            \textbf{First look, then leap} \\
            It is better not to “pollute” the nanopub template space. First look or ask for existing templates.
        \end{callouttip}
        \vspace{1em}
        \begin{calloutnote}
            We are \textbf{not} going to create a new template in this workshop, but we will briefly show you how it’s done.
        \end{calloutnote}
    }
    
%%%%%%%%%%%%%%%%%%%%%%%%%%%%%%%%%%%%%%%%%%%%%%%%%%%%%%%%%%%%%%%%%%%%%%%%%%%%%%%%  
    \frame{
        \frametitle{Modifying a Template II}
        \begin{itemize}
            \item<+-> Start by selecting a suitable (similar) template and click the little blue caret \\
            \includegraphics[width=0.4\linewidth]{new_template_01.png}
            \item<+-> Then select \textbf{“edit as derived nanopublication”} \\
            \includegraphics[width=0.6\linewidth]{new_template_02.png}
        \end{itemize}
    }

%%%%%%%%%%%%%%%%%%%%%%%%%%%%%%%%%%%%%%%%%%%%%%%%%%%%%%%%%%%%%%%%%%%%%%%%%%%%%%%%    
    \frame{
        \frametitle{Modifying a Template III}
        \begin{callouttip}
            \textbf{Start small} \\
            Creating a new assertion template can be difficult. Start with an easy (or small) one.
        \end{callouttip}
        \vspace{1em}
        \begin{callouttip}
            \textbf{Ask for input!} \\
            If you need substantial changes, ask for help. (In the end we will show some resources.)
        \end{callouttip}
    }

%%%%%%%%%%%%%%%%%%%%%%%%%%%%%%%%%%%%%%%%%%%%%%%%%%%%%%%%%%%%%%%%%%%%%%%%%%%%%%%%
    \frame{
        \frametitle{What can you do with Nanopubs?}
        Construct knowledge-graphs for your paper!
        \vspace{1em}
        \centering
        \includegraphics[width=0.8\linewidth]{Snakemake_Metadata.png} \\
        \vspace{0.5em}
        \footnotesize Created with \href{https://github.com/snakemake/snakemake-report-plugin-nanopub/}{the Snakemake reporter plugin for posting workflow metadata}.
        \vspace{1em}
        \begin{callouttip}
            \textbf{Citing Nanopubs} \\
            Remember: nanopubs are unique persistent identifiers, worldwide accessible.
        \end{callouttip}
    }

%%%%%%%%%%%%%%%%%%%%%%%%%%%%%%%%%%%%%%%%%%%%%%%%%%%%%%%%%%%%%%%%%%%%%%%%%%%%%%%%    
    \frame{
        \frametitle{What can you do with Nanopubs? II}
        More \& more Citation Ontologies become available in Scholia:
        \vspace{1em}
        \centering
        \includegraphics[width=0.8\linewidth]{CiTO_Scholia.png} \\
        \vspace{1em}
        \begin{callouttip}
            \textbf{Nanopub has a Template for your CiTOs} \\
            \href{https://w3id.org/np/RAX_4tWTyjFpO6nz63s14ucuejd64t2mK3IBlkwZ7jjLo}{Declare CiTOs with this template}
        \end{callouttip}
    }

%%%%%%%%%%%%%%%%%%%%%%%%%%%%%%%%%%%%%%%%%%%%%%%%%%%%%%%%%%%%%%%%%%%%%%%%%%%%%%%%    
    \frame{
        \frametitle{And now? Where to go from here?}
        This workshop \textit{merely} gave an overview. There is so much more:
        \begin{itemize}
            \item<+-> \textbf{Community}
            \begin{itemize}
                \item<+-> \href{https://nanopub.net/sessions/}{Nano Sessions} – every last Tuesday of a month; reports from practitioners; ask questions; \\
                \href{https://nanodash.knowledgepixels.com/space?id=https://w3id.org/spaces/nanopub/nanosessions}{register here with a nanopub}
                \item<+-> \href{https://nanodash.knowledgepixels.com/space?id=https://w3id.org/spaces/fenac}{FENAC} – the “Fearless Early Nanopub Adopters Club”; sign-up per nanopub; members get prioritized feature requests and a free one-on-one call with a core developer.
            \end{itemize}
            \item<+-> \textbf{Resources}
            \begin{itemize}
                \item<+-> \href{https://nanopub.net/docs/tutorials/}{Official tutorial selection}
                \item<+-> \href{https://nanopub.net/docs/}{Documentation}
            \end{itemize}
            \item<+-> \textbf{Software with an API for automated publication}
            \begin{itemize}
                \item<+-> \href{https://github.com/Nanopublication/nanopub-java}{Java library}
                \item<+-> \href{https://github.com/Nanopublication/nanopub-py/}{Python library}
            \end{itemize}
            \item<+-> \textbf{Software that uses the API}
            \begin{itemize}
                \item<+-> \href{https://github.com/snakemake/snakemake-report-plugin-nanopub/}{Snakemake reporter plugin for posting workflow metadata}
            \end{itemize}
        \end{itemize}
    }
    
